\subsection{Roadmap}
\begin{frame}{Roadmap: how the rest of today addresses the requirements}
  \framesubtitle{One requirement at a time — then an \textit{in practice} act applying them all}

  \setlength{\tabcolsep}{5pt}
  \renewcommand{\arraystretch}{1.3}
  \begin{center}
    \footnotesize
    \begin{tabular}{@{}p{4.8cm}p{7.0cm}@{}}
      \toprule
      \textbf{Requirement} & \textbf{Addressed by}\\
      \midrule
      \textbf{Req. \ref{req:universal}} --- serve multiple use cases &
        Kinds of information models: separating \emph{vocabulary} from \emph{constraint};
        later: profiles as use-case-specific L2 models\\
      \textbf{Req. \ref{req:technology}} --- technology-agnostic &
        MOF: metamodels, expressivity, choosing a \emph{pivot model},
        metamodel transformations; later: generated artifacts (Avro, JSON Schema, SQL)\\
      \textbf{Req. \ref{req:consistent}} --- no logical inconsistencies &
        Upper ontologies \& ontology alignment: spotting and fixing semantic conflicts (the
        \texttt{cim:Outage} example)\\
      \bottomrule
    \end{tabular}
  \end{center}

  \vspace{0.1cm}
  \begin{block}{The goal}
    A practical guide to how to meet all these requirements. Looking at how Netbeheer Nederland addresses them. And finally looking back at openADR.
  \end{block}
\end{frame}
